\chapter{Dirac vs Copenhagen: Two Foundational Positions}\label{ch:dirac}
\index{Dirac, P. A. M.}\index{Copenhagen interpretation}\index{Born interpretation}\index{Schr\"odinger, E.}\index{configuration space}\index{StdQM, standard quantum mechanics}\index{wave function}\index{parameter-free, Level 1}

% Material to be added later from prel book (realquantum2.tex):
%   - chapters/something.tex, somethingfresh.tex (Solvay 1927, the Step)
%   - chapters/postulates.tex (Pauli on exclusion, the official story)
%   - chapters/timeline.tex (1925--1930 chronology)

\begin{quote}
\itshape There is a difference between a shaky or not sharply focused photograph and a photograph of clouds and fogbanks.\\
\normalfont\hfill --- E.~Schr\"odinger, \emph{The Present Situation in Quantum Mechanics}, 1935.
\end{quote}
\bigskip

The Prologue (page~\pageref{ch:prologue}) sketched the development of quantum mechanics from Planck in 1900 to Dirac in 1929 and identified the decisive choice point of 1926--1929: whether to extend Schr\"odinger's 1926 wave equation to the $3N$-dimensional configuration space, as Born, Heisenberg, and Dirac proposed, or to keep the wave function on physical 3D space, as Schr\"odinger himself had originally envisaged. The configuration-space extension won, became standard quantum mechanics (StdQM), and produced the chemistry-specific machinery that has organised the field ever since. This chapter elaborates that choice point, frames it as two foundational positions --- Dirac's chemistry-as-applied-physics and the Copenhagen interpretation --- and previews the alternative the rest of the book develops.

\section{Two quotations}\label{sec:two-quotations}

The book opens with two quotations from the founders of quantum mechanics. They were written within three years of each other, by physicists who knew each other's work intimately, and they sketch two different futures for the new theory. One of those futures became orthodoxy. The other --- the one this book takes up --- was abandoned almost immediately and has not been seriously pursued in mainstream physics since.

The first is from \textbf{Paul Dirac, 1929}, on the relation between physics and chemistry:

\begin{quote}
The underlying physical laws necessary for the mathematical theory of [\ldots] the whole of chemistry are thus completely known, [and] the difficulty is only that the exact application of these laws leads to equations much too complicated to be soluble.
\end{quote}

The second is from \textbf{Erwin Schr\"odinger, in a letter to H.~A.\ Lorentz dated 6 June 1926}, sketching how he originally intended the wave function to be interpreted for systems with more than one electron:

\begin{quote}
If we now have to deal with $N$ particles, then $\Psi(x_1, \ldots, x_N)$ is a function of $N$ variables $x_1, \ldots, x_N$ over $N$ 3D spaces $R_1, \ldots, R_N$. Now first let $R_1$ be identified with the real space $\mathbb{R}^3$ and integrate over $R_2, \ldots, R_N$. Second, identify $R_2$ with the real space and integrate over $R_1, R_3, \ldots, R_N$ and so on. The $N$ individual results are to be added after they have been multiplied by certain constants which characterise the particles. \emph{I consider the result to be the electric charge density in real space.}
\end{quote}

The two positions are not complementary; they are in tension. Dirac states that chemistry reduces to known physical laws; the only obstacle is computational --- but he treats the configuration-space form of the equations as fixed and concludes that the only relief is approximation. Schr\"odinger, by contrast, sketches the real-space alternative: the $N$-electron wave function projects to an actual charge density in physical $\mathbb{R}^3$, recovering the direct physical meaning that $\psi^2$ has for the hydrogen atom. On Schr\"odinger's picture the equation Dirac calls ``too complicated'' is too complicated because it was written on the wrong space; rewriting it on physical $\mathbb{R}^3$ removes the exponential cost without weakening the physics. The choice between these two positions, taken in 1926--1929, decided which of the two futures became standard quantum mechanics.

Schr\"odinger himself remained sceptical of the probabilistic interpretation throughout his life --- the cat-paper ``fogbanks'' line that opens this chapter is one record of his unease, and others are scattered as epigraphs throughout the book. The \textbf{real-space-density interpretation was abandoned during the late 1920s} under the dominance of what came to be called the Copenhagen interpretation. In that reading, $\Psi$ on $\mathbb{R}^{3N}$ became a probability amplitude rather than a physical charge density, and the multi-dimensional configuration space became the proper home of the wave function. The resulting framework --- which we will call Standard QM, or \textbf{StdQM}, throughout this book --- inherits the computational obstacle Dirac identified: solving an equation on $3N$-dimensional configuration space. The chemistry-specific approximations introduced over the following decades to make StdQM tractable have, ever since, made chemistry look more autonomous from physics than Dirac's position would suggest.

\section{The configuration-space step}\label{sec:configuration-space}

The move that decided which of the two futures became orthodoxy was the extension of the wave function from physical $\mathbb{R}^3$ to a $3N$-dimensional configuration space. It is worth being precise about what this step did and what it cost, because the rest of this book is organised around \emph{not taking it}.

For the hydrogen atom Schr\"odinger's 1926 wave function $\psi(\mathbf{x})$ is a real-valued function on physical 3D space. Its square $\psi^2(\mathbf{x})$ has the direct physical meaning of an electronic charge density: at each point $\mathbf{x}$ of real space, $\psi^2$ tells you how much of the electron's charge sits there. There is no probability and no observer; there is a continuous charge cloud bound to a proton, with shape and total energy determined by minimisation of a definite functional. The hydrogen problem was solved, and solved decisively, in this picture.

For two or more electrons the picture had to be extended, and there were two ways to do it. Schr\"odinger's letter to Lorentz sketches the first: keep the wave function on physical $\mathbb{R}^3$, with one such function per particle, and combine them into a single charge density on real space by adding the squared amplitudes (with appropriate species charges). The second way, suggested by Born and developed by Heisenberg --- and adopted by Schr\"odinger himself in his published work --- writes a single wave function $\Psi(\mathbf{x}_1, \ldots, \mathbf{x}_N)$ on the $3N$-dimensional product space and reads $|\Psi|^2$ as a probability density for finding electron 1 at $\mathbf{x}_1$, electron 2 at $\mathbf{x}_2$, and so on, simultaneously.

The two extensions are not equivalent and do not represent the same physics. The first stays in physical 3D and gives up the convenience of writing the $N$-body problem as a single eigenvalue equation. The second buys mathematical compactness at the price of leaving physical space and accepting a probabilistic interpretation that the hydrogen problem did not require. The second won, for a mixture of mathematical and sociological reasons, and the real-space picture was set aside.

\paragraph{The sociological side: the Forman thesis.}\index{Forman thesis}\index{Weimar Republic and quantum mechanics}
The phrase ``sociological reasons'' deserves a sentence of unpacking because the historiography on this point is not marginal. Paul Forman's 1971 thesis, ``Weimar Culture, Causality, and Quantum Theory'', proposed that the acausal, probabilistic reinterpretation of the wave function that took hold in 1926--1929 was adopted with unusual speed and with an unusual absence of internal resistance partly because acausality was already congenial to the intellectual atmosphere of Weimar Germany --- a culture in which Spenglerian fatalism, Lebensphilosophie, and a broad rejection of mechanistic causality were dominant strands in the wider intellectual life. The configuration-space-plus-probability step was, on Forman's reading, not only a mathematical move but also a cultural one, and its near-immediate acceptance owed something to its consonance with the surrounding intellectual mood as well as to its technical merits. The historical-sociology literature continues to debate the Forman thesis, but the broader point --- that the rapid consolidation of the Copenhagen reading was not driven by mathematical necessity alone --- is widely accepted, and a less culturally-specific causal-realist alternative might well have remained in serious contention had the timing been different. The framework developed in this book is, in part, an attempt to revisit that alternative with a century of distance from the Weimar moment.

The cost of the configuration-space step was felt almost immediately:

\begin{itemize}
\item \emph{Computational.} Representing $\Psi$ on a $K$-point grid in physical space costs $O(K)$ in memory; representing it on a grid in $3N$-dimensional configuration space costs $O(K^N)$. A 100-grid-point calculation that is trivial for one electron becomes infeasible for ten. Every quantum-chemistry method developed since 1929 --- Hartree--Fock, configuration interaction, coupled cluster, density functional theory --- can be read as an attempt to manage this exponential cost, by truncating the variational space, by replacing the full wave function with a single Slater determinant, or by passing entirely to a one-body density.

\item \emph{Interpretational.} Once $\Psi$ lives on $\mathbb{R}^{3N}$, $|\Psi|^2$ no longer has the direct physical meaning that $\psi^2$ had for hydrogen. It is a joint probability density on a space that is not physical space. The recovery of an actual charge density in $\mathbb{R}^3$ requires integrating out $N-1$ coordinates and is conventionally described in probabilistic language even when the underlying configuration is stationary. The probabilistic reading was not necessary, as Schr\"odinger's letter to Lorentz makes clear: a real-space charge-density reading is available, takes the form $\rho(\mathbf{x}) = \sum_i \psi_i^2(\mathbf{x})$, and gives a deterministic field on physical space. The probabilistic reading was adopted nonetheless, and the foundational debate it generated has continued for a century not because the question is genuinely open but because the wrong formulation was institutionalised.

\item \emph{Conceptual.} The exponential cost is not just a numerical inconvenience; it changes what kind of object the wave function is. $\Psi$ on $\mathbb{R}^{3N}$ is not a field on physical space the way the electromagnetic field is; it is a function on a much larger abstract space. The intuition that ``the wave function is a thing in the world'' --- straightforward for the hydrogen $\psi$ --- becomes, for $N>1$, a contested matter that physicists have debated for nearly a century.
\end{itemize}

It is conventional to present the configuration-space step as a forced move, the only mathematically respectable way to extend the one-electron equation to many electrons. Dirac's 1929 sentence cuts against that convention: the underlying physics is known, the equations are merely \emph{too complicated}. \emph{Too complicated} is a property of the equations one chose to write, not of the physics. If a different choice yields equations that are not too complicated, the underlying physics has not changed; only the formulation has.

\section{RealQM: the formulation Schr\"odinger should have reached}

The formulation developed in this book, which we call \textbf{RealQM}, returns to the originally proposed real-3D-space picture. The $N$-electron wave function is written as
\begin{equation}\label{eq:Psi-sum-intro}
\Psi(\mathbf{x}) \;=\; \sum_{i=1}^{N} \psi_i(\mathbf{x}),
\qquad \mathbf{x} \in \mathbb{R}^3,
\end{equation}
a sum of one-electron wave functions $\psi_i$, each defined on its own region $\Omega_i$ of physical 3D space, with $\{\Omega_i\}$ a non-overlapping partition of $\mathbb{R}^3$. Each $\psi_i^2$ is a unit charge density: the actual electronic charge of electron $i$ in its own territory. The total wave function $\Psi$ inherits the dimensionality of physical space; it is a field on $\mathbb{R}^3$, not on $\mathbb{R}^{3N}$.

The energy is the standard Coulomb energy --- kinetic, kernel--electron, electron--electron --- evaluated as an integral over $\mathbb{R}^3$. The ground state is the configuration of $\{\psi_i\}$, the partition $\{\Omega_i\}$, and the nuclear positions $\{\mathbf{R}_a\}$ that minimises this energy. Molecular geometry, bonding, and dynamics emerge from the same variational principle.

The configuration-space step is not taken at all. The computational obstacle Dirac identified is dissolved by working in physical 3D throughout, at the structural cost of replacing antisymmetry of an $N$-body wave function with strict spatial non-overlap of one-body wave functions. The two requirements are not equivalent. They produce different mathematical objects with different scaling, different stability proofs, and different relationships to chemistry.

In this respect RealQM is a direct realisation of \textbf{Dirac's chemistry-as-applied-physics position} executed in the real-3D-space picture Schr\"odinger originally envisaged. RealQM takes Dirac's program seriously by retaining only the underlying physics --- Coulomb interactions between electronic wave functions and nuclei --- and discarding the chemistry-specific machinery (basis sets, exchange--correlation functionals, fitted force fields) that StdQM had to introduce in order to make the configuration-space equations tractable. Chemistry then re-emerges, bottom-up, from the variational principle alone.

This reformulation is not merely a numerical recasting of the standard problem. It is a different mathematical object with different scaling: complexity grows with the number of mesh points in 3D, not with the number of basis functions in $N \times 3D$. As a consequence, the entire framework can be implemented in roughly 8000 lines of JavaScript and WGSL compute shaders, runs in a browser on a consumer laptop GPU, and supports interactive simulation of systems with hundreds of explicit electrons in real time.

\section{Hierarchy of reductions: a preview}

A theoretical framework is worth taking seriously to the extent that it works across scales. RealQM is organised as a hierarchy of reductions in which each level is derived from the previous one by a specific operation, and each level is checked against observation before being used as the foundation of the next.

\textbf{Level 1} is parameter-free: an atom is described as a nucleus of charge $Z$ surrounded by $N$ one-electron wave functions $\psi_i$ on non-overlapping shell-shaped domains, with the configuration determined by minimum-energy packing. Level 1 reproduces observed atom total energies for elements Li through Rn to approximately 1\% with no fitted constants.

\textbf{Level 2} spherically homogenises inner shells, keeping the shell occupancy intact but replacing the angular structure of $\sum_{i \in \text{shell}} \psi_i^2$ with a radial density of matching total charge. The bond-relevant chemistry is unaffected because inner shells contribute to the effective Coulomb potential seen by valence electrons, not to bond formation directly.

\textbf{Level 3} replaces the nucleus and inner shells together by a \emph{pseudo-kernel} characterised by an effective charge $Z_{\rm kernel}$ and a softening radius $r_c$, leaving only the outermost valence wave functions explicit. The kernel parameters are determined by comparison with the Level-2 results, in the same architectural spirit as pseudopotentials and frozen-core methods in standard quantum chemistry but with a parameter-free ab initio bottom rather than empirical fits.

\textbf{Level 4} assembles Level-3 atoms into molecules and molecular complexes.

The hierarchy continues upward as a vision toward biology. \textbf{Level 5} treats whole amino-acid residues, rather than individual atoms, as the units of representation --- the natural scale for protein folding and protein--ligand docking. \textbf{Level 6} treats whole proteins, lipid bilayers, and small subcellular structures as Level-6 objects in a cellular context --- the natural scale for protein--protein recognition, membrane systems, and supramolecular assembly. Levels 5 and 6 are partially in use already (residue-level docking, peptide folding) and partially aspirational; Chapter~\ref{ch:hierarchy} describes their status in detail. Each level is the input to the next; each level is calibrated against the previous one. The validation chapters in Part~III work through Levels 1--4 quantitatively; Level 0, the atomic nucleus in proton--electron form, is presented in Chapter~\ref{ch:nucleus} at the same exploratory stage as Levels 5 and 6 occupy at the upper end.

\section{The collaboration as part of the contribution}

A second contribution of this book concerns how the implementation behind it was produced. The mathematical reformulation is the work of one author. The validation suite, the WebGPU compute shaders, the systematic kernel-architecture sweeps reported in the validation chapters, and the curated public Gallery were developed through extended interactive collaboration with \textbf{Claude}, a large-language-model AI assistant produced by Anthropic.

The pattern of the collaboration was iterative and recorded. The human proposed mathematical and chemical questions; the AI translated them into runnable simulations, identified bugs, ran systematic parameter sweeps, and challenged claims that did not survive scrutiny; the human evaluated the results and corrected the AI's interpretations when they slipped. The history of corrections --- claims made, retracted, refined, restored --- is itself part of the public record, encoded in the Gallery's ``Assessment by Claude'' card.

We take this case to be a worth-reporting example of how a single mathematical mind, working with an AI engineering partner in real time, can produce research-grade interactive scientific software at a scale otherwise requiring a programming team. The collaboration is not the science, but it is part of how the science came into existence, and Chapter~\ref{ch:collab} reports it as such.

\section{Plan of the book}

Before the six numbered parts, two unnumbered front pieces set the wider perspective. The \textbf{Foreword} places the volume in the Body and Soul programme: a 3D parameter-free continuum-mechanics formulation of quantum mechanics paired with a runnable implementation, presented as the realisation of Leibniz's three-and-a-half-century-old programme of automating computation, made tractable in the present moment by the working partnership of a single mathematical mind with capable AI assistance. The \textbf{Prologue} (\emph{Schr\"odinger 1926--1952}) opens with the strange/rational diagnosis --- standard quantum mechanics is described as strange by its own practitioners while being praised as the most empirically successful theory in history, and the diagnosis reads the strangeness as the price of the $\mathbb{R}^3 \to \mathbb{R}^{3N}$ configuration-space step that this book reverses --- then sets out the history from Planck through Dirac with Schr\"odinger's 1927 Solvay communication and his 1952 Dublin Colloquium as the conceptual brackets. The book is organised in six parts.

\textbf{Part I (Foundations)} states the reformulation. Chapter~\ref{ch:equation} writes down the multiphase 3D Schr\"odinger equation, derives its Euler--Lagrange equations, introduces the Bernoulli free-boundary condition that closes the system, and develops the spatial-non-overlap replacement of antisymmetry together with the geometric reading of Pauli's Zweideutigkeit, the AIM-basin and Leibniz-monad correspondences, and the connection to classical electromagnetism through charge density. Chapter~\ref{ch:hierarchy} sets out the four levels of reduction in detail, including the recovery of Lewis's 1916 cubic atom as the Level-1 ground state of the L shell.

\textbf{Part II (Numerical method)} describes how the equations are solved. Chapter~\ref{ch:relaxation} presents the three coupled time-evolution equations --- imaginary-time propagation of the $\psi_i$, parabolic-relaxation Poisson updates of the per-electron Hartree potentials, and level-set evolution of the free boundary --- that constitute the entire solver. Chapter~\ref{ch:webgpu} describes how those three equations are mapped onto WebGPU compute shaders that run interactively in a browser.

\textbf{Part III (Validation)} works through the hierarchy of reductions against observation. Chapter~\ref{ch:atoms} reports Level-1 atomic energies from Li to Rn. Chapter~\ref{ch:h2} reports the H$_2$ covalent bond and its He limit. Chapter~\ref{ch:hydrides} reports closed-shell hydride atomization energies across periodic-table groups. Chapter~\ref{ch:dimers} reports geometric agreement against the S66 benchmark of non-covalent dimers, together with reactive chemistry. Chapter~\ref{ch:folding} reports interactive protein folding. Chapter~\ref{ch:condensed} reports a three-system condensed-phase calibration (NaCl crystal melt, CO$_2$ dry-ice cluster, solid N$_2$). Chapter~\ref{ch:materials} places the condensed-phase results in their materials-science context.

\textbf{Part IV (Foundations revisited)} returns to foundational questions in light of the validation. Chapter~\ref{ch:som} shows that stability of matter --- the lower bound $E \ge -CNZ^3$ on the ground-state energy of a system of $N$ atoms --- follows directly from the Hardy inequality applied to non-overlapping electronic territories, in contrast to the heroic Dyson--Lenard--Lieb--Thirring proof required in standard quantum mechanics. Chapter~\ref{ch:nucleus} reports an analogous model of the atomic nucleus, demonstrating that the framework extends beyond chemistry to nuclear-scale physics.

\textbf{Part V (Matter and Light)} extends the framework into matter--light interaction. Chapter~\ref{ch:sm-coupling} sets out the Schr\"odinger--Maxwell coupling in which the RealQM charge density and current source the Maxwell field directly on physical 3D space. Chapter~\ref{ch:excited} reports excited states and atomic radiation. Chapter~\ref{ch:molrad} reports molecular radiation with worked examples for H$_2$O and CO$_2$. Chapter~\ref{ch:blackbody} reports a Rayleigh--Jeans-truncated derivation of the black-body spectrum with Stefan--Boltzmann and Wien-displacement scaling preserved. This part is the most speculative reach of the book; the chapters in it are at a different stage of validation from the chemistry chapters of Part~III.

\textbf{Part VI (Connections)} places the work in its broader intellectual and institutional context. Chapter~\ref{ch:thermodynamics} states the deterministic-continuum reading of thermodynamics that the matter--light chapters require, parallel to \emph{Computational Thermodynamics} (Body and Soul Vol~5). Chapter~\ref{ch:chemvphys} returns to the chemistry-vs-physics debate in light of what RealQM does and does not do, and contains the institutional-diagnosis cluster of the book: a discussion of why the Hilbert-space framework does not distinguish atomic physics from classical continuum mechanics; the von Neumann split between the 1932 QM formalism and the 1945 stored-program computer architecture; the engagement of mathematicians with StdQM that started in 1928--1936 with Hilbert--von Neumann--Weyl, was abandoned mid-stream, and was never picked up by a comparable later figure; the displacement of foundational work into the discipline of philosophy of physics housed outside physics departments, with Bohr as the founding philosopher; Einstein's lifelong objection from EPR to Bell; and the doubling-down pattern by which QED and then string theory were built on top of unresolved foundational difficulties rather than addressing them. Chapter~\ref{ch:collab} discusses the human--AI collaboration in detail. Chapter~\ref{ch:limits} states honest limits of the approach, including the double-slit experiment (clean RealQM story), entanglement (outside present scope), and the Stern--Gerlach two-fold splitting (a different two-valuedness from the Pauli Zweideutigkeit already explained in Part~I). The Appendix lists the source repositories, the public Gallery, the hardware requirements, and the \emph{Leibniz World of Math} development site that has documented the framework's development over more than a decade.

\section{What this book is, and is not}

This book is a long-form expansion of the Foundations-of-Physics article \emph{Quantum Mechanics as Multiphase 3D Continuum Mechanics}~\cite{Johnson2026FoP}. The article is dense and compact; the book is meant to be read in sequence by someone who wants to understand both the mathematical content and the route by which it was arrived at. Each chapter expands one section of the article, with room left for historical material from the author's earlier 2017 draft and for newer results that did not fit into the article.

This book is also the long-planned \textbf{quantum-mechanics volume of the Body and Soul programme}~\cite{BodyAndSoulI,BodyAndSoulII,BodyAndSoulIII,BodyAndSoulIV,BodyAndSoulV}, the author's two-decade constructive--computational reformulation of applied mathematics. Body and Soul carried calculus, differential equations, fluid mechanics, and other classical subjects through to runnable algorithms paired with the underlying mathematics; the quantum-mechanics volume was on the agenda from the beginning but waited until the implementation became feasible for a single mathematical mind paired with an AI engineering partner. Chapter~\ref{ch:collab} discusses that precursor relationship in detail. Readers familiar with the earlier volumes will recognise the methodological continuity: every claim is paired with a simulation, every reduction is calibrated against observation, every limit is identified empirically.

\paragraph{No figures: read the book in parallel with the GitHub Gallery.}
This book \textbf{does not contain any figures from the many code examples} that underlie the validation in Part~III and the rest of the framework. This is a deliberate choice. The Gallery hosts the live interactive simulations --- density slices, 3D molecular views, force-arrow displays, real-time relaxation, parameter sweeps --- that produce every quantitative claim in the book. Static images of those simulations omit precisely what makes the framework distinctive: that the reader can change a parameter and watch the chemistry respond.

\textbf{The reader is therefore strongly recommended to read this book in parallel with the GitHub Gallery, with all the code examples open to run and modify.} A chapter typically describes \emph{what} is computed and \emph{why} it agrees with observation; the corresponding Gallery card lets the reader \emph{see} it, vary it, and check the agreement directly. The repository, the Gallery URL, the per-chapter pointers, and the hardware requirements are listed in the Appendix. A modern browser and a few minutes of setup are all that is needed.

This book is \emph{not} a textbook on standard quantum chemistry. It does not derive Hartree--Fock, density functional theory, or coupled cluster from first principles, and does not try to replicate the existing pedagogical literature on those methods. References to StdQM are made for comparison, not for instruction.

This book \emph{is} an argument that standard quantum mechanics, in its configuration-space formulation, was a wrong turn taken in 1926--1929, and that the framework it generated has consumed nearly a century of foundational work without resolving the difficulties it created. The position taken throughout is direct: \emph{here is the formulation Schr\"odinger should have arrived at; here is what it predicts; here is how it agrees with experiment; and here is what it costs to compute --- which is a laptop and a browser}. The standard formulation is not presented as a parallel approach to be weighed against this one; it is presented as the apparatus that grew up around the wrong choice, with results that are right where the underlying physics is right (the Coulomb force law, the original real-space Schr\"odinger equation) and incoherent where its foundational story has had to be filled in (the probabilistic interpretation, antisymmetry as primitive, the measurement problem, hidden variables, many-worlds). The reader is asked to consider, in light of what this book demonstrates, whether the configuration-space step is still defensible as the foundation of quantum mechanics, or whether the time has come to return to the 1926 picture and start from there.

\paragraph{The submission record.}
For factual completeness, the route that led from the manuscript to this book is recorded here. The article \emph{Quantum Mechanics as Multiphase 3D Continuum Mechanics}~\cite{Johnson2026FoP} was submitted to \emph{Foundations of Physics} on 2 May 2026 and was declined by editorial decision in May 2026 \textbf{without referee report and without stated motivation}. The same manuscript was then submitted to arXiv (physics.atom-ph / physics.chem-ph), where it was refused for posting; arXiv's moderation routinely declines manuscripts that propose foundational reformulations of quantum mechanics, regardless of the technical level of the content or the institutional affiliation of the author. The book-length manuscript --- the volume the reader is holding --- was offered to Springer through Martin Peters, the editor of \emph{Body and Soul Vols I--IV}, with no response; in the editorial context of Springer's physics imprint, that silence reads as a soft decline that preserves the personal relationship while avoiding the institutional cost of either advocacy or formal rejection. None of these channels engaged with the manuscript's technical content. The framework, the validation reported in Part~III, and the demonstration that the calculation fits on a laptop and runs in a browser have not been disputed: they have been declined to be heard.

\paragraph{A case of suppressed scientific debate.}
The pattern is not unique to this manuscript. It is the same pattern that affected \textbf{David Bohm}'s 1952 papers on a deterministic interpretation of quantum mechanics (declined by the canonical physics journals until \emph{Physical Review} accepted them only after sustained pressure, and dismissed institutionally for decades thereafter), \textbf{Hugh Everett}'s 1957 relative-state thesis (effectively buried for fifteen years before DeWitt's 1972 revival), and \textbf{Louis de Broglie}'s pilot-wave programme (sidelined for nearly thirty years after the 1927 Solvay Conference). The mechanism in each case was not refutation but \emph{non-engagement}: editors of the canonical venues declined to send the work out for review, with the result that the standard formulation continued to be presented in the literature and in education as the only available formulation regardless of what alternative formulations were actually in circulation. The intellectual cost of this pattern, by now well documented in the history-of-science literature~\cite{Cushing94, Freire15, Beller99}, has been a century in which the foundational questions raised by Einstein, Schr\"odinger, de Broglie, and Bohm have been treated as settled philosophy rather than as open physics --- not because their technical content was disposed of, but because the technical content was institutionally prevented from reaching the audience that could have engaged with it. The Dublin Coda of the Prologue, in which Schr\"odinger himself states the position in 1952, is the founding sample in this lineage; the present submission record is one further data point in the same century-long pattern.

\paragraph{What this book does about it.}
This volume, and the public Gallery that accompanies it, are an attempt to step around the pattern rather than to argue against it from within. The standard journal-and-arXiv route is, on the available evidence, closed to foundational reformulations of quantum mechanics as a matter of editorial policy, regardless of the technical merit of any specific submission. What remains open is direct publication of the long-form record, with running code that the reader can execute, modify, and check independently of any editorial gatekeeping. The reader is therefore invited to do the one thing the standard channels have so far declined to do: read the manuscript, run the code, examine the results, and decide for themselves whether the framework reaches across the scales it claims to reach. The institutional verdict on whether RealQM is or is not the right way to think about quantum mechanics is, at the time the present volume is being put into the reader's hands, not yet rendered --- because the institutions have so far declined to render one.

% --- End of Chapter 1 ---
